
\medskip
\noindent\textbf{0.2.}
\phantomsection
\label{sga3:VIIB:0.2}
Let \(A\) be a commutative pseudocompact ring. A
\emph{pseudocompact \(A\)-module} \(M\) is a separated, complete
topological \(A\)-module having a basis of neighborhoods of \(0\)
consisting of submodules \(M'\) for which \(M/M'\) has finite length
over \(A\).

If \(M\) and \(N\) are pseudocompact \(A\)-modules, a
\emph{morphism} from \(M\) to \(N\) is, by definition, a continuous
\(A\)-linear map. We write
\[
  \operatorname{Hom}_c(M,N)
\]
for the group of these morphisms.

\medskip
\noindent\SGAInlineRefTarget{sga3:vii:proposition:0-2-b}\textbf{Proposition 0.2.B.}\footnote{The editors have stated
the results of this paragraph separately in the following proposition,
and have then indicated the steps of the proof; cf.\ \textup{[CA]},
\S~IV.3, or \textup{[DG70]}, \S~V.2.}
\phantomsection
\label{sga3:VIIB:0.2.B}
{\itshape
\begin{enumerate}
  \item[\textup{(i)}]
  The pseudocompact \(A\)-modules form an abelian category, denoted
  \(\operatorname{PC}(A)\). In particular, for every morphism
  \(f:M\to N\), the submodule \(\operatorname{Im}(f)\) is complete and
  hence closed in \(N\).

  \item[\textup{(ii)}]
  The pseudocompact \(A\)-modules of finite length, whose topology is
  therefore discrete, form a system of cogenerators of
  \(\operatorname{PC}(A)\).

  \item[\textup{\SGAInlineRefTarget{sga3:vii:proposition:0-2-b:item:iii}(iii)}]
  Infinite products and filtered projective limits are exact; in other
  words, \(\operatorname{PC}(A)\) satisfies axiom
  \((\mathrm{AB}5^{*})\).
\end{enumerate}}

For the reader's convenience, we briefly indicate the steps of the
proof. First we use the following lemma (\textup{[CA]}, \S~IV.3,
Lemma~1; for the proof, see \textup{[BEns]}, III, \S~\SGARefLink{sga3:VIIA:7.4}{7.4}, Theorem~1
and Example~2).

\medskip
\noindent\textbf{Lemma 0.2.C.}
\phantomsection
\label{sga3:VIIB:0.2.C}
\textit{Let \(B\) be a ring, let \(I\) be a filtered ordered set, and
let \((M_i)\) and \((N_i)\) be two projective systems of left
\(B\)-modules indexed by \(I\). Let
\[
  (s_i):(M_i)\longrightarrow(N_i)
\]
be a morphism of projective systems such that \(s_i\) is surjective
and has Artinian kernel for every \(i\). Then the map}
\[
  \varprojlim_i s_i:
  \varprojlim_i M_i
  \longrightarrow
  \varprojlim_i N_i
\]
\textit{is surjective.}

\medskip
\noindent\textbf{Corollary 0.2.D}
\textit{(\textup{[CA]}, \S~IV.3, Propositions~10 and~11).}
\phantomsection
\label{sga3:VIIB:0.2.D}
{\itshape Let \(M\) be a pseudocompact \(A\)-module.
\begin{enumerate}
  \item[\textup{(i)}]
  If \(K\) is a closed submodule of \(M\), then \(M/K\), endowed with
  the quotient topology, is a pseudocompact \(A\)-module.

  \item[\textup{(ii)}]
  Let \((M_i)\) be a filtered decreasing family of closed submodules
  of \(M\).
  \begin{enumerate}
    \item[\textup{\SGAInlineRefTarget{sga3:VIIB:0.2.D:item:a}(a)}]
    The canonical map
    \[
      M\longrightarrow\varprojlim_i M/M_i
    \]
    is surjective and has kernel \(\bigcap_i M_i\).

    \item[\textup{\SGAInlineRefTarget{sga3:VIIB:0.2.D:item:b}(b)}]
    For every closed submodule \(N\) of \(M\), one has
    \[
      N+\bigcap_i M_i
      =
      \bigcap_i(N+M_i).
    \]
  \end{enumerate}
\end{enumerate}}

